\documentclass[11pt]{article}

% ------------------------------------------------------------------
% MiniChess AI -- Mini Project 2 report
% Compile: pdflatex report.tex  (uses only standard packages)
% ------------------------------------------------------------------
\usepackage[utf8]{inputenc}
\usepackage[margin=1in]{geometry}
\usepackage{amsmath,amssymb}
\usepackage{booktabs}
\usepackage{xcolor}
\usepackage{listings}
\usepackage{enumitem}
\usepackage[hidelinks]{hyperref}

\definecolor{codebg}{rgb}{0.96,0.96,0.96}
\definecolor{kw}{rgb}{0.0,0.0,0.6}
\definecolor{cmt}{rgb}{0.0,0.45,0.0}
\lstset{
  basicstyle=\ttfamily\small,
  backgroundcolor=\color{codebg},
  keywordstyle=\color{kw}\bfseries,
  commentstyle=\color{cmt}\itshape,
  breaklines=true,
  frame=single,
  columns=fullflexible,
  showstringspaces=false,
  language=C++,
  numbers=none,
}
\newcommand{\code}[1]{\texttt{#1}}

\title{\textbf{MiniChess AI -- Design, Search, and a Neural-Network Evaluation Study}\\
\large Mini Project 2}
\author{Student ID: \texttt{114000110}}
\date{\today}

\begin{document}
\maketitle

\begin{abstract}
This report documents a strong AI for the 6$\times$5 ``MiniChess'' variant. The
engine combines a classical alpha-beta / Principal Variation Search (PVS) with a
full modern search stack (transposition table, killer/counter/history move
ordering, iterative deepening, aspiration windows, null-move pruning, late-move
reductions, futility pruning, quiescence search) and a \emph{hybrid evaluation}
that blends a hand-crafted heuristic with a small \emph{NNUE} (Efficiently
Updatable Neural Network) trained from scratch on self-play data. Beyond the
working engine, the report presents an empirical study: porting Stockfish's NNUE
architecture ideas (SFNNv1--v14) to MiniChess, int16 quantization, texel tuning,
and a search-allocation optimization. The central, repeatedly-confirmed finding
is that \textbf{on a tiny board, search depth dominates evaluation
sophistication}, and that \textbf{training-loss improvements do not imply
playing-strength improvements}. The final engine beats every provided baseline
(weak/strong minimax, alpha-beta, PVS) with both colours and the visible boss.
\end{abstract}

\tableofcontents

% ==================================================================
\section{Overview and architecture}
% ==================================================================
The program speaks \textbf{UBGI} (a UCI-superset protocol) over stdin/stdout.
The code is split into game-agnostic and game-specific layers:

\begin{itemize}[nosep]
  \item \code{src/ubgi/ubgi.cpp} -- protocol loop and the \emph{iterative
        deepening} driver (calls the search at depth $1,2,3,\dots$, honours
        \code{go movetime}, emits \code{info}/\code{bestmove}).
  \item \code{src/games/minichess/state.cpp} -- board representation, move
        generation (naive + bitboard), Zobrist hashing, terminal detection.
  \item \code{src/policy/submission.cpp} -- the engine we developed: PVS search,
        all pruning/ordering heuristics, the hand-crafted evaluation, and the
        NNUE blend.
  \item \code{src/policy/nnue.hpp}, \code{nnue\_weights.hpp} -- the neural
        evaluator and its embedded weights.
  \item \code{tools/nnue\_train.cpp} -- a self-contained C++ self-play data
        generator and trainer (no external ML dependencies).
\end{itemize}

\paragraph{Rules and representation.} MiniChess is chess on a 6-row $\times$
5-column board with a \emph{king-capture} win model: a side wins the instant it
can capture the enemy king. There is no castling and pawns only promote to a
queen. The board is stored as \code{char board[2][6][5]} (one plane per player;
values $1{=}\text{P},2{=}\text{R},3{=}\text{N},4{=}\text{B},5{=}\text{Q},6{=}\text{K}$).
Draws occur by repetition and a step cap. A useful symmetry: Black's starting
array is exactly White's rotated $180^\circ$, which the NNUE exploits
(Section~\ref{sec:nnue}).

% ==================================================================
\section{State value function (evaluation)}
% ==================================================================
The evaluation returns an integer score in centipawn-like units \emph{from the
side-to-move's perspective} (so the search is a pure negamax). The final score
is a \textbf{blend} of a hand-crafted term and a neural term.

\subsection{Hand-crafted evaluation}
\code{custom\_evaluate()} computes, for each side,
\begin{equation}
S_{\text{hand}} = \big(\text{material} + \text{PST} + \text{king tropism}\big)_{\text{self}}
  - (\cdots)_{\text{opp}} + \text{pawn terms} + w_m\cdot\text{mobility},
\end{equation}
with the following components:
\begin{itemize}[nosep]
  \item \textbf{Material}: $\{P,R,N,B,Q,K\}=\{100,500,320,330,900,10000\}$.
  \item \textbf{Piece-square tables (PST)}: a per-piece $6\times5$ table adding
        positional value; the table is read from the side-to-move's orientation.
  \item \textbf{King tropism}: a bonus for own pieces being close to the enemy
        king (encourages attacking the king, which directly ends the game).
  \item \textbf{Pawn structure}: penalties for doubled and isolated pawns.
  \item \textbf{Passed pawns}: a bonus scaling with proximity to promotion
        (max near the promotion rank), added because promotion is decisive.
  \item \textbf{Mobility}: $w_m=5$ times the number of legal moves available to
        the side to move.
\end{itemize}
Terminal positions short-circuit: a win returns $P_{\max}$, a draw returns $0$.

\subsection{NNUE evaluation}\label{sec:nnueeval}
A small neural network (Section~\ref{sec:nnue}) produces a score from the same
side-to-move perspective. It is blended with the hand-crafted score:
\begin{equation}
S = \frac{b\cdot S_{\text{nnue}} + (100-b)\cdot S_{\text{hand}}}{100},
\qquad b = \texttt{NNUE\_BLEND} = 30.
\end{equation}
NNUE is on by default and self-contained; \code{-DNO\_NNUE} disables it and
\code{-DNNUE\_BLEND=$n$} overrides the mix. \textbf{Why a blend, and why only
30\%?} Empirically (Section~\ref{sec:blend}) pure NNUE is \emph{weaker} than the
hand-crafted eval on this game, but a light 30\% admixture is meaningfully
stronger than either component alone: the net captures outcome-correlated
patterns the heuristic misses, while the heuristic supplies tactical robustness
and is cheap to compute.

% ==================================================================
\section{Tree search}
% ==================================================================
\subsection{From minimax to negamax}
Minimax assumes the opponent minimises our score. Because MiniChess strictly
alternates turns and the evaluation is symmetric (side-to-move relative), we use
the \emph{negamax} form: $\text{value}(n) = \max_{m}\,-\,\text{value}(\text{child}(n,m))$.

\begin{lstlisting}[language=C++,caption={Negamax with alpha-beta (concept)}]
int negamax(node, depth, alpha, beta):
    if depth == 0 or terminal: return evaluate(node)   // (-> quiescence)
    for each move m:
        score = -negamax(child(m), depth-1, -beta, -alpha)
        if score >= beta: return beta          // beta cutoff (fail-high)
        if score > alpha: alpha = score
    return alpha
\end{lstlisting}

\subsection{Alpha-beta pruning}
$\alpha$ is the best score the side to move is already assured of; $\beta$ is the
best the opponent is assured of. When $\alpha \ge \beta$ the remaining moves at a
node cannot affect the result (the opponent would avoid this line), so the
branch is pruned. Alpha-beta returns the \emph{same} value as minimax but visits
far fewer nodes; with good move ordering it approaches the theoretical
$O(b^{d/2})$ instead of $O(b^d)$.

\subsection{Principal Variation Search (PVS)}
PVS assumes the first (best-ordered) move is best: it searches that move with a
full window $[\alpha,\beta]$ and every later move with a \emph{null window}
$[\alpha,\alpha{+}1]$ (a cheap test of ``is this move worse than the best so
far?''). If a null-window search unexpectedly fails high, it is re-searched with
the full window. This is implemented in both \code{minimax.cpp} (the reference)
and \code{submission.cpp} (the tuned engine).

\subsection{Quiescence search}
At depth 0 the engine does not evaluate immediately; it runs a
\emph{quiescence} search over ``noisy'' moves (captures and promotions) until the
position is quiet, preventing the \emph{horizon effect} (mis-evaluating a
position in the middle of a capture sequence). It uses a stand-pat cutoff, MVV-LVA
ordering, and \emph{delta pruning} (skip captures that cannot raise $\alpha$ even
with a generous margin).

\subsection{Transposition table (TT)}
A $2^{22}\approx4{\cdot}10^6$-entry hash table caches search results keyed by the
Zobrist hash. Each \code{TTEntry} ($\approx$10\,bytes) stores a 32-bit key check,
score, depth, a bound flag (exact / lower / upper), a 16-bit packed best move
(\code{CompactMove}), and a generation byte for ageing. The TT (i)~returns a
cached score when a position is reached again at sufficient depth and (ii)~feeds
the \emph{TT move} to move ordering, which is the single biggest ordering win.
Mate scores are stored ply-adjusted so they remain correct across transpositions.

\subsection{Move ordering}
Good ordering makes alpha-beta/PVS effective. Moves are tried in this priority:
\begin{enumerate}[nosep]
  \item \textbf{TT move} (the best move found previously at this position).
  \item \textbf{Captures} by \textbf{MVV-LVA} (Most Valuable Victim $-$ Least
        Valuable Attacker: $16\cdot V_{\text{victim}} - V_{\text{attacker}}$).
  \item \textbf{Promotions}.
  \item \textbf{Killer moves}: two quiet moves per ply that caused a cutoff.
  \item \textbf{Counter move}: the quiet move that previously refuted the
        opponent's last move.
  \item \textbf{History heuristic}: quiet moves scored by how often they have
        caused cutoffs (indexed by from/to square).
\end{enumerate}

\subsection{Selectivity: pruning, reductions, extensions}
On top of plain alpha-beta the engine searches \emph{selectively}:
\begin{itemize}[nosep]
  \item \textbf{Null-move pruning (NMP)}: give the opponent a free move; if the
        position is still $\ge\beta$ at reduced depth ($R=\min(d{-}1,\,3{+}[d\!\ge\!7])$,
        clamped $\ge2$), prune. Disabled in check and at PV nodes.
  \item \textbf{Late-move reductions (LMR)}: late, quiet, non-killer moves are
        searched at reduced depth (reduction $1$--$4$ growing with move index and
        depth, less at PV nodes); if a reduced search beats $\alpha$ it is
        re-searched at full depth.
  \item \textbf{Reverse futility / static null move}: at low depth, if the static
        eval exceeds $\beta$ by $120\cdot d$, return early.
  \item \textbf{Futility pruning}: at low depth, skip quiet moves when
        $\text{static}+150\cdot d \le \alpha$.
  \item \textbf{Late-move pruning (LMP)}: at very low depth, skip quiet moves
        beyond a move-count threshold.
  \item \textbf{Mate-distance pruning}: tighten the window using known mate
        bounds so faster mates are preferred.
  \item \textbf{Check extension}: search one ply deeper when in check.
  \item \textbf{Internal iterative deepening (IID)}: at PV nodes lacking a TT
        move, do a shallow search first to obtain a good first move.
\end{itemize}

\subsection{Iterative deepening, aspiration windows, time management}
The driver searches depth $1,2,\dots$, reusing the previous depth's results
(TT + PV) for ordering. From the previous score it opens a narrow
\emph{aspiration window} ($\pm 50$); on a fail-high/low it widens exponentially
and re-searches. Time is checked every 2048 nodes; the search aborts shortly
before the per-move budget using a margin of $\min(300,\,\text{movetime}/25)$\,ms.
A ``half-time'' rule avoids starting a new depth that almost certainly cannot
finish. \textbf{The competition limit is a strict 2 seconds per move}; the engine
adapts to whatever \code{movetime} the runner passes.

% ==================================================================
\section{NNUE: a neural evaluation trained from scratch}\label{sec:nnue}
% ==================================================================
No pre-trained weights or ML libraries were used; the entire pipeline is in C++.

\subsection{Feature encoding (side-to-move relative)}
The board is canonicalised to the side to move so the same weights serve both
colours. White keeps orientation; \textbf{Black is rotated $180^\circ$}
($r\!\to\!5{-}r,\ c\!\to\!4{-}c$), which is exact because Black's start array is
White's rotated $180^\circ$. Each occupied square contributes one binary feature:
\begin{equation}
\text{feature} = \underbrace{\text{base}}_{0\ \text{own},\ 180\ \text{enemy}}
 + (\text{piece}-1)\cdot 30 + \text{sq},\quad
\text{sq}=\begin{cases} r\cdot5+c & \text{White}\\ (5{-}r)\cdot5+(4{-}c) & \text{Black}\end{cases}
\end{equation}
giving $2\times6\times30 = 360$ features (own/enemy $\times$ 6 piece types $\times$
30 squares). The encoding is implemented identically in the trainer and the
engine -- a mismatch would silently corrupt the evaluation.

\subsection{Network and training pipeline}
The shipped network is a small fully-connected net $360 \to 32 \to 1$ with ReLU
and a $\tanh$ output (scaled by $600$ to centipawns). Training (\code{nnue\_train.cpp}):
\begin{enumerate}[nosep]
  \item \textbf{Self-play data}: $3000$ games are played (random openings of
        $1$--$8$ plies plus $12\%$ random moves for diversity, otherwise a
        depth-4 search picks moves), yielding $\approx 98\text{k}$ positions.
  \item \textbf{Labels (WDL)}: every position is labelled with the eventual game
        result from \emph{its} side-to-move's perspective ($+1$ win, $0$ draw,
        $-1$ loss). Outcome labels carry information beyond the current eval,
        which is the whole point.
  \item \textbf{Training}: Adam, mini-batch MSE against the $\tanh$ output,
        $60$ epochs (final MSE $\approx 0.56$).
  \item \textbf{Embedding}: weights are written to \code{nnue\_weights.hpp} as
        plain C++ arrays, so the engine is self-contained (no runtime model
        file) and compiles with the standard command.
\end{enumerate}

\subsection{Inference: sparse, per-leaf, float}
The board has $\le 20$ pieces, so the accumulator is rebuilt \emph{sparsely from
scratch} at every leaf -- summing the $\le 20$ active feature rows of $W_1$ --
rather than incrementally. This avoids make/unmake accumulator bugs and is cheap
enough that NPS stays high ($\sim$1.4\,M nodes/s). Output:
\begin{equation}
\text{cp} = \mathrm{round}\Big(\tanh\big(b_2 + \textstyle\sum_h W_2[h]\cdot\mathrm{relu}(b_1[h]+\!\!\sum_{f\,\text{active}}\!\!W_1[f][h])\big)\cdot 600\Big).
\end{equation}

\subsection{Blend tuning}\label{sec:blend}
We measured several blend ratios with \emph{diverse openings} (see
Section~\ref{sec:method}). Pure NNUE ($b{=}100$) lost to a $b{=}50$ blend
$0$/$8$; higher NNUE weight was monotonically worse ($b{=}70$ scored 25\% vs
$b{=}50$). The chosen $b{=}30$ beats the pure hand-crafted engine
\textbf{62.5\%} and the boss \textbf{87.5\%} (diverse openings), and sweeps the
boss \textbf{8/8} from the fixed start at 2\,s.

% ==================================================================
\section{Architecture study: porting Stockfish NNUE (SFNNv1--v14)}
% ==================================================================
We investigated whether Stockfish's NNUE architecture evolution transfers to
MiniChess. This section is the most likely demo discussion point.

\subsection{What changed across SFNN versions}
\begin{table}[h]\centering\small
\begin{tabular}{@{}lll@{}}
\toprule
Version & Key architectural idea & Transfers to MiniChess? \\
\midrule
SFNNv1 & HalfKP features; \emph{two-perspective} accumulator; & dual-persp / clipped-ReLU: yes; \\
       & clipped ReLU; deeper head; int quantization & HalfKP: too many features \\
SFNNv4 & HalfKAv2 (king included as a piece) & partial (we include the king) \\
SFNNv5 & HalfKAv2\_hm (horizontal mirroring) & no -- board is not L/R symmetric \\
SFNNv6 & L1 width doubled; \emph{piece-count output buckets} & buckets: cheap, plausible \\
v7--v14& L1 scaling $1024\!\to\!3072$, more data & scaling only \\
\bottomrule
\end{tabular}
\end{table}

\subsection{Dual-perspective accumulator (SFNNv1 core)}
We implemented the defining NNUE idea: encode the position from \emph{both} the
side-to-move's and the opponent's viewpoint through a \emph{shared} feature
transformer, concatenate the two accumulators, and pass through a deeper head
($360 \to 64\,[\times2] \to 32 \to 1$) with clipped ReLU. Trained on pure WDL it
reached a much lower MSE ($0.237$ vs $0.56$) -- yet it \textbf{played worse}: the
two accumulators plus the bigger head cut NPS $\sim$3$\times$
($1.26\text{M}\to350\text{k}$), costing $\sim$1.5--2 plies of depth. Result:
25\% vs hand-crafted, 80\% vs boss (down from 62.5\%/87.5\%).

\subsection{int16 quantization (auto-vectorized)}
To recover the lost speed we quantized to int16 (SFNNv1's speed trick): weights
scaled by $Q_A{=}512$ (clipped-ReLU unit) and $Q_W{=}256$ (layer weights), with
\emph{int32} accumulation laid out contiguously so the compiler auto-vectorizes
under \code{-march=native} (float-vs-int parity $\approx$2.5\,cp). This raised the
dual-perspective net to 502\,k\,nps and 45\%/85\% -- a clear recovery, but
\emph{still below} the simple single-perspective net. Quantizing the
single-perspective net instead gave \emph{no} speedup (its float eval is already
cheap and int16 widening cancels the lane-count gain) and the clipped-ReLU
retraining slightly regressed it (47.5\% vs hand).

\subsection{Conclusion of the study}
\textbf{Stockfish's NNUE architecture is sound but does not transfer to
MiniChess.} On full chess, NNUE's eval cost is amortised by enormous search
trees, the hand-crafted alternative is far weaker, and int8 SIMD + incremental
updates keep NNUE cheap. On a $6\times5$ board none of that holds: a cheap eval
plus more search depth wins, so the structural sophistication (and even its
quantization speed trick) is a net loss. The shipped engine therefore uses the
\emph{single-perspective float} net. (All experiments are preserved on git
branches as documented negative results.)

% ==================================================================
\section{Search-speed optimization: a ply-indexed node pool}
% ==================================================================
Profiling pointed at per-node memory churn: the search did \code{new
State}+\code{delete} (and a fresh move-list allocation) for every node. We added
\code{State::apply\_into} / \code{null\_into}, which write the child position
\emph{in place} into a caller-owned \code{State} taken from a ply-indexed pool
(\code{g\_node\_pool}), reusing each slot's move-list buffer. Depth-first search
guarantees \code{pool[child\_ply]} is free to overwrite once a subtree returns,
and a parent never reads a child slot while recursing -- so the search tree is
\emph{provably identical} (verified: same node counts and scores) while two heap
allocations per node disappear. Net effect: \textbf{$\sim$12\% higher NPS
($1.26\text{M}\to1.41\text{M}$)}, validated at 10/12 vs the boss. This is the one
optimization (of the three attempted) that shipped.

% ==================================================================
\section{Negative results (also useful at the demo)}
% ==================================================================
\begin{itemize}[nosep]
  \item \textbf{Texel-tuning the material values} to fit WDL improved the static
        loss by only $0.4\%$ and \emph{lost 37.5\%} head-to-head: static-eval
        loss is a poor proxy for strength, and material values also drive move
        ordering / delta pruning, so reweighting hurt search efficiency.
  \item \textbf{Opening book}: with a strict \emph{per-move} 2\,s limit (no game
        clock) a book cannot bank time and the engine already plays the fixed
        opening optimally, so a book would only echo the engine's own move. Not
        shipped (it would be a no-op).
  \item \textbf{Search-score teacher for NNUE} (labelling positions with the
        engine's own deep-search score instead of game outcome) reached a very
        low MSE but regressed in play -- the net merely learned to imitate the
        existing eval.
\end{itemize}

% ==================================================================
\section{Results and methodology}\label{sec:method}
% ==================================================================
\paragraph{Measuring strength correctly.} Games from the fixed start position are
nearly deterministic (only timing jitter varies), so a ``20-game match'' is
effectively one game repeated and cannot resolve small differences. We therefore
added \emph{random openings} (\code{--opening-plies N --seed S}) to the match
runner; all comparative numbers below use diverse openings.

\begin{table}[h]\centering\small
\begin{tabular}{@{}lll@{}}
\toprule
Opponent & Result (final engine) & Notes \\
\midrule
Weak MiniMax (baseline)   & 4/4    & both colours \\
Strong MiniMax (baseline) & 4/4    & both colours \\
Alpha-Beta (baseline)     & win    & weaker than the PVS boss we beat 8/8 \\
PVS (visible boss)        & 8/8    & 2\,s, fixed start \\
PVS boss, diverse openings& 87.5\% & 20 games, harder test \\
Hand-crafted-only engine  & 62.5\% & isolates the NNUE blend's value \\
\bottomrule
\end{tabular}
\caption{Final engine (single-perspective NNUE blend, node pool) vs.\ opponents.}
\end{table}

\paragraph{Constraints respected.} Standard library only (NNUE weights are
embedded C++ arrays); no manual SIMD/inline-asm (only compiler auto-vectorization
via \code{-march=native}); the engine's policy code is single-threaded;
TT $\approx$48\,MB, well under the 4\,GB limit.

% ==================================================================
\section{Key points to remember for the demo}
% ==================================================================
\begin{enumerate}[nosep]
  \item \textbf{Search}: negamax + alpha-beta + PVS + quiescence, plus TT,
        killer/counter/history ordering, IID, aspiration windows, NMP, LMR,
        reverse-futility/futility/late-move pruning, mate-distance pruning, and
        check extensions. Be ready to explain \emph{why} each prunes safely.
  \item \textbf{Evaluation}: material + PST + king tropism + pawn structure +
        passed pawns + mobility, \emph{blended} with a from-scratch NNUE at 30\%.
  \item \textbf{NNUE}: 360-feature side-to-move-relative encoding ($180^\circ$
        flip for Black), $360\!\to\!32\!\to\!1$, trained by C++ self-play on WDL
        labels, weights embedded in a header, sparse per-leaf inference.
  \item \textbf{Headline insight}: \emph{on a small board, depth beats eval
        sophistication}; SOTA chess NNUE architecture and its quantization trick
        do \emph{not} transfer (measured, not assumed).
  \item \textbf{Second insight}: \emph{lower training loss $\ne$ stronger play}
        -- shown three times (score-teacher NNUE, dual-perspective NNUE, texel
        tuning).
  \item \textbf{Methodology}: use random openings to measure strength; a clean
        speedup (node pool) must be a provable non-regression (identical node
        counts) before shipping.
\end{enumerate}

\paragraph{Build.}
\begin{lstlisting}[language=bash]
g++ --std=c++2a -Wall -Wextra -Wpedantic -g -O3 -march=native -static \
  -Isrc/games/minichess -Isrc/state -Isrc \
  -o build/minichess-ubgi.exe \
  src/games/minichess/state.cpp src/policy/minimax.cpp \
  src/policy/random.cpp src/policy/submission.cpp src/ubgi/ubgi.cpp
\end{lstlisting}
NNUE is on by default and self-contained; \code{-DNO\_NNUE} disables it and
\code{-DNNUE\_BLEND=$n$} changes the mix.

\end{document}
